% Beamer's default font size is 11 points. It is possible to set the
% default font size to any of 8, 9, 10, 11, 12, 14, 17, 20
\documentclass[13pt]{beamer}

\usetheme{Warsaw}
%\usecolortheme{dolphin}
\usecolortheme{seahorse}
%\useinnertheme{circles}
\usefonttheme{professionalfonts}
%\usefonttheme{structurebold}
%\usefonttheme{default}

\setbeamertemplate{navigation symbols}{}
 
\usepackage[english]{babel}
\usepackage[latin1]{inputenc}
\usepackage{times}
\usepackage[T1]{fontenc}
\usepackage{fancybox}

% These things can change:
\newcommand{\baseurl}{http://cell.vtt.fi/latex}
\urlstyle{sf}
\newcommand{\plus}{\makebox[2em]{+}}
\newcommand{\minus}{\makebox[2em]{-}}
\newcommand{\gray}[1]{%
\textcolor{gray}{#1}
}

% In beamer, defining new commands inside the slides fails.
% Definitions here:

\begin{document}

\logo{\includegraphics[width=1.4cm]{img/vttplain}\quad}
\title{\scalebox{1.5}{\textrm{\LaTeX}} Course 2011}
\subtitle{Part 6: Companion tools}
\author{Arho Virkki}
\institute{\textsc{VTT Technical Research Centre of Finland}}
\date{\scalebox{0.8}{\textcolor{gray}{}}}

\begin{frame}[fragile]
  \titlepage
\end{frame}

\begin{frame}[fragile]\frametitle{The need for extra tools}
As seen earlier, LaTeX is a capable system for typesetting. For productive
work, we still need 
\begin{itemize}
\item citation management system,
\item advanced drawing tools, 
\item perhaps some more exotic LaTeX add-on packages, and
\item a proper version control system (e.g. cvs, svn or 
git)\footnote{These tools are not discussed in this course. Nevertheless, LaTeX
users get their full power, as the manuscripts are text -- similar to software
source code.\\\mbox{}}
\end{itemize}
\end{frame}


\begin{frame}[fragile]\frametitle{BibTeX}

BibTeX is the de facto standard for LaTeX citatioin management
(written by Oren Patashnik and Leslie Lamport in 1985). It has
\begin{itemize}
\item An uncomplicated structure:
\begin{itemize}
  \item the database can be modified by hand and
  \item it is easy to write programs that can handle BibTeX files.
\end{itemize}
\item A lot of BibTeX tools already exists, including
\begin{itemize}
  \item \texttt{bibtool}
  \item \texttt{bibtex2html}
  \item \dots
\end{itemize}
\end{itemize}
\end{frame}


\begin{frame}[fragile]\frametitle{BibTeX\dots}
BibTeX-database is just a plain text file. Example:
\begin{verbatim}
@Article{gro67,
  author  = {Fred S. Grodins and ...},
  title   = {Mathematical analysis and ...},
  journal = {Journal of Applied Physiology},
  volume  = {22},
  number  = {2},
  pages   = {260-276},
  year    = {1967},
  url     = {Grodins - Maadsotrcs.pdf}
}
\end{verbatim}
\end{frame}


\begin{frame}[fragile]\frametitle{BibTeX\dots}
\begin{itemize}
\item The record begins with the type, for example \verb|@article|,
  \verb|@book| or \verb|@techreport|.
\item The first string is the key for the citation, which can be used in the
text as \verb|\cite{lyhenne}|.
\item All fields are key-value pairs, like \verb|pages={260-276}|.
\item All fields are separated by commas.
\end{itemize}
The format of the record is hence
\begin{verbatim}
@tyyppi{field, field, ..., field}
\end{verbatim}
\end{frame}


\begin{frame}[fragile]\frametitle{BibTeX\dots}
\begin{itemize}
\item The authors are introduced in the order \emph{first\_name last\_name} and
\item the names are separated with keyword \texttt{and}.
\end{itemize}
For example, 
\begin{verbatim}
J. Doe and O. Normalverbraucher and 
M. Meikalainen
\end{verbatim}\medskip
\begin{flushleft}
The style of the Bibliography is defined by a separate style sheet. This can
usually be downloaded from the home page of the journal the article is being
submitted.
\end{flushleft}
\end{frame}


\begin{frame}[fragile]\frametitle{BibTeX\dots}
The style sheet can e.g. define that the title begins with a capital letter, and
all the following words are written in lower case. In this case, for example, 
\begin{verbatim}
Stability of the Human Respiratory 
Control System I: Analysis of a 
Two-Dimensional Delay State-Space Model
\end{verbatim}
becomes
\begin{verbatim}
Stability of the human respiratory 
control system i: analysis of a 
two-dimensional delay state-space model
\end{verbatim}
which was not necessarily the original intention.
\end{frame}


\begin{frame}[fragile]\frametitle{BibTeX\dots}
If we need to fix the presentation of a word, that part of the text can be
written inside extra curly braces \{ \}
\begin{verbatim}
... System {I}: {A}nalysis of a ...
\end{verbatim}
which makes these parts immutable.
\end{frame}


\begin{frame}[fragile]\frametitle{BibTeX and natbib}
Using the add-on package
\begin{verbatim}
\usepackage{natbib}
\end{verbatim}
we get two additional commands in addition to the basic \verb|\cite{}|:
\begin{itemize}
\item \verb|\citet{}| --- the \emph{textual} and
\item \verb|\citep{}| --- the \emph{parenthetical citation}.
\end{itemize}
that enable taking the citations as part of the text.
\end{frame}


\begin{frame}[fragile]\frametitle{BibTeX and natbib\dots}
Example:
\begin{verbatim}
\citet{Foo88} showed that... furthermore ...
in the last decade \citep{Bar93,Baz96}.
\end{verbatim}
produces
\begin{quotation}
\noindent
Fooman et al. (1988) showed that... furthermore ...
in the last decade (Barnos et al., 1993; Bazel et al., 1996).
\end{quotation}
depending of the chosen style.
\end{frame}



\begin{frame}[fragile]\frametitle{BibTeX usage}
Example: \medskip
If BibTeX database is stored in the file \texttt{citations.bib},
it can be cited in the text using the command \verb|\cite{key}|, provided that
the file has been added to the list of citation files:
\medskip
\begin{verbatim}
\bibliographystyle{apalike}
\bibliography{citations}
\end{verbatim}
where \texttt{apalike.bst} is one of the default style sheets available.
\end{frame}


\begin{frame}[fragile]\frametitle{BibTeXin usage\dots}
In case that the document (here \verb|text.tex|) is compiled for the first time,
one needs to run all the commands
\begin{verbatim}
latex text 
bibtex text 
latex text 
latex text 
\end{verbatim}
until all is done.
\end{frame}

\begin{frame}[fragile]\frametitle{Style sheet tricks}
One of our very first examples had the command
\begin{verbatim}
@Article{gro67,
  ...,
  url     = {Grodins - Maadsotrcs.pdf}
}
\end{verbatim}
\begin{flushleft}
where \texttt{url} refers to the file name in the local machine. This, of
course, has no meaning for other users, and the field is best to be removed from
the public version of the file.
\end{flushleft}
\end{frame}


\begin{frame}[fragile]\frametitle{Style sheet tricks\dots} 
Because the \texttt{.bib} files are just ordinary text, removing a field is
easy.\medskip

Example: (\texttt{sed} = strean editor)
\begin{verbatim}
sed --in-place '/url/ d' citations.bib
\end{verbatim}
The same can also be done with a proper Python or R program.
\end{frame}


\begin{frame}[fragile]\frametitle{Beamer}
Beamer is an additional package for writing slides.
\begin{itemize}
\item You write ordinary \LaTeX and
\item Prosper makes to output to look like it was produced using ``the power
point method''.
\item --- except for the formulas, which will still be typeset with the  
\LaTeX quality.
\item The style definitions of the current slide set are 
\begin{verbatim}
\usetheme{Warsaw}
\usecolortheme{seahorse}
\usefonttheme{professionalfonts}
\end{verbatim}
\end{itemize}
\end{frame}


\begin{frame}[fragile]\frametitle{Beamer\ldots}
Beamer comes installed by defauld with LiveTeX, 
and compiling the slides (in \verb|slides.tex|) is straightforward:
\begin{verbatim}
pdflatex slides
\end{verbatim}
\begin{itemize}
\item You had better to choose a program that automatically notices the update 
\texttt{pdf}-file.
\item Under linux,  \texttt{evince} automatically notices any change in the
file. For Mac OS X, the editor usually incorporates a proper previewer.
\end{itemize}
\end{frame}

\begin{frame}[fragile]\frametitle{Beamer example}
The following commands tell LaTeX that the document is a slide set (and sets
some extra preferences)
\begin{verbatim}
\documentclass[13pt]{beamer}
\usetheme{Warsaw}
\usecolortheme{seahorse}
\usefonttheme{professionalfonts}
\setbeamertemplate{navigation symbols}{}

\begin{document}
...
\end{verbatim}
\end{frame}


\begin{frame}[fragile]\frametitle{Beamer example...}
Each slide begins and ends with 
\begin{verbatim}
  \begin{frame}[fragile]
  \frametitle{The slide title}
  This goes to the slide
  ...
  \end{frame}
\end{verbatim}
The extra argument \texttt{[fragile]} is not always necessary, but it enables
using the \texttt{verbatim} environment.\footnote{Which is heavily used e.g.
for these slides.\\\mbox{}}
\end{frame}


\begin{frame}[fragile]\frametitle{Beamer example...}
\normalsize
These slides are actually bad. Good slides have as little text as
possible, and\medskip
\begin{itemize}
\bfseries
\item the text is large enough and it
\item has lots of free space around it.
\end{itemize}\bigskip

In addition,\medskip
\begin{itemize}
\sl
\item each slide contains at most 4 key elements, and 
\item the audience is not expected to read and listen at the same time.
\end{itemize}
\end{frame}


\end{document}
